\begin{tikzpicture}[
  every node/.style={font=\footnotesize\sffamily},
  participant/.style={draw=TFGBlue, fill=TFGLightBlue, rounded corners=2pt,
    minimum width=2.25cm, minimum height=0.65cm, align=center},
  lifeline/.style={dashed, draw=TFGMuted},
  message/.style={-{Latex[length=2mm]}, semithick, draw=TFGAccent},
  return/.style={-{Latex[length=2mm]}, dashed, draw=TFGMuted},
  message label/.style={fill=white, inner xsep=2pt, inner ysep=1pt,
    text=black, align=center}
]
  \node[participant] (user) at (0,0) {Miembro};
  \node[participant] (web) at (3.8,0) {Aplicación web};
  \node[participant] (api) at (7.6,0) {Gestor documental};
  \node[participant] (storage) at (11.4,0) {Almacenamiento};
  \node[participant] (db) at (15.2,0) {Persistencia};
  \foreach \x in {0,3.8,7.6,11.4,15.2} \draw[lifeline] (\x,-0.35) -- (\x,-7.05);
  \draw[message] (0,-0.75) -- node[message label, above]{selecciona un PDF} (3.8,-0.75);
  \draw[message] (3.8,-1.45) -- node[message label, above]{envía archivo y datos asociados} (7.6,-1.45);
  \draw[message] (7.6,-2.15)
    -- node[message label, above]{valida formato y tamaño} ++(0.75,0)
    -- ++(0,-0.55) -- (7.6,-2.7);
  \draw[message] (7.6,-3.35) -- node[message label, above]{almacena el contenido} (11.4,-3.35);
  \draw[return] (11.4,-4.05) -- node[message label, above]{referencia del objeto} (7.6,-4.05);
  \draw[message] (7.6,-4.75) -- node[message label, above]{registra los metadatos} (15.2,-4.75);
  \draw[return] (15.2,-5.45) -- node[message label, above]{registro confirmado} (7.6,-5.45);
  \draw[return] (7.6,-6.15) -- node[message label, above]{documento disponible} (3.8,-6.15);
  \draw[return] (3.8,-6.75) -- node[message label, above]{confirma la incorporación} (0,-6.75);
\end{tikzpicture}
